\chapter*{Abstract}
Electroencephalography (EEG) records the brain's electrical activity and may help researchers study differences between people with Alzheimer's disease and healthy individuals. However, a classifier that gives a correct answer for one part of a recording may change its answer later, while a model that never changes its answer can still be consistently wrong. This thesis first examines these problems using subject-level EEG classification and then studies a custom decoder built on a pretrained EEG model. The first set of experiments used a traditional support vector machine to evaluate Alzheimer's disease, frontotemporal dementia and healthy-control recordings across time. It found that first-minute balanced accuracy was about 73\% for both binary comparisons, but predictions changed during the five-minute assessment for many participants. We then developed EEGraft's gated dual-scale residual fusion decoder, which combines features from a frozen CBraMod encoder with features learned from the EEG signal. On 52 development participants, this decoder reached 77.06\% balanced accuracy, compared with 76.61\% for XGBoost. On a separate locked group of 13 participants, the decoder reached 70.24\%, while XGBoost reached 77.38\%. Artificial noise and channel-removal tests also showed that prediction stability and accuracy are different properties. These results support careful subject-level evaluation and show that a custom decoder can be competitive without consistently outperforming a classical model. The findings are limited to retrospective data and do not establish a clinical diagnosis or predict future Alzheimer's disease.
\vspace{0.5cm}
\noindent\textbf{Keywords:} Alzheimer's disease; EEG; CBraMod; dual-scale residual fusion; subject-level evaluation; temporal stability; robustness; reliability.
